\noindent The remaining algebra tables collect the recovery formulas, product decompositions, and order-dependent component expansions used in later calculations.
\begin{componenttable}[tab:components-paravector-recovery]{Z}
\componentrow{scalar}{a+ib=(Z+Z^{\mathrm{R}*})/2}{}
\componentrow{vector}{(\mathbf A+i\mathbf B)\cdot\sigv=(Z-Z^{\mathrm{R}*})/2}{}
\componentrow{real part}{a+\mathbf A\cdot\sigv=(Z+Z^{\mathrm R})/2}{}
\componentrow{imaginary part}{b+\mathbf B\cdot\sigv=(Z-Z^{\mathrm R})/(2i)}{}
\componentrow{real scalar}{a=(Z+Z^*+Z^{\mathrm R}+Z^{\mathrm{R}*})/4}{}
\componentrow{imaginary scalar}{b=(Z-Z^*-Z^{\mathrm R}+Z^{\mathrm{R}*})/(4i)}{}
\componentrow{real vector}{\mathbf A\cdot\sigv=(Z-Z^*+Z^{\mathrm R}-Z^{\mathrm{R}*})/4}{}
\componentrow{imaginary vector}{\mathbf B\cdot\sigv=(Z+Z^*-Z^{\mathrm R}-Z^{\mathrm{R}*})/(4i)}{}
\end{componenttable}

\noindent The following two tables collect the full component expansions for the generic product \(ZZ'\) and the special case \(Z^2\).
\begin{componenttable}[tab:components-paravector-product]{Z Z'}
\componentrow{real scalar}{aa'-bb'+\mathbf A\cdot\mathbf A'-\mathbf B\cdot\mathbf B'}{}
\componentrow{imaginary scalar}{ab'+ba'+\mathbf A\cdot\mathbf B'+\mathbf B\cdot\mathbf A'}{}
\componentrow{real vector}{a\mathbf A'-b\mathbf B'+a'\mathbf A-b'\mathbf B-\mathbf A\times\mathbf B'-\mathbf B\times\mathbf A'}{}
\componentrow{imaginary vector}{a\mathbf B'+b\mathbf A'+a'\mathbf B+b'\mathbf A+\mathbf A\times\mathbf A'-\mathbf B\times\mathbf B'}{}
\end{componenttable}

\begin{componenttable}[tab:components-paravector-square]{Z^{2}}
\componentrow{real scalar}{a^{2}-b^{2}+\mathbf A^{2}-\mathbf B^{2}}{}
\componentrow{imaginary scalar}{2ab+2\mathbf A\cdot\mathbf B}{}
\componentrow{real vector}{2(a\mathbf A-b\mathbf B)}{}
\componentrow{imaginary vector}{2(a\mathbf B+b\mathbf A)}{}
\end{componenttable}

\noindent The following two tables make the order dependence of \(ZZ^{*}\) and \(Z^{*}Z\) explicit.
\begin{componenttable}[tab:components-zzstar]{Z Z^{*}}
\componentrow{real scalar}{a^{2}+b^{2}-\mathbf A^{2}-\mathbf B^{2}}{}
\componentrow{imaginary scalar}{0}{}
\componentrow{real vector}{-2\,\mathbf A\times\mathbf B}{\(\leftarrow\) differs from \(Z^{*}Z\)}
\componentrow{imaginary vector}{2(a\mathbf B-b\mathbf A)}{}
\end{componenttable}

\begin{componenttable}[tab:components-zstarz]{Z^{*} Z}
\componentrow{real scalar}{a^{2}+b^{2}-\mathbf A^{2}-\mathbf B^{2}}{}
\componentrow{imaginary scalar}{0}{}
\componentrow{real vector}{2\,\mathbf A\times\mathbf B}{\(\leftarrow\) differs from \(ZZ^{*}\)}
\componentrow{imaginary vector}{2(a\mathbf B-b\mathbf A)}{}
\end{componenttable}

\noindent The following two tables play the same role for products involving the reverse operator.
\begin{componenttable}[tab:components-zzrev]{Z Z^{\mathrm R}}
\componentrow{real scalar}{a^{2}+b^{2}+\mathbf A^{2}+\mathbf B^{2}}{}
\componentrow{imaginary scalar}{0}{}
\componentrow{real vector}{2(a\mathbf A+b\mathbf B+\mathbf A\times\mathbf B)}{\(\leftarrow\) differs from \(Z^{\mathrm R}Z\)}
\componentrow{imaginary vector}{\mathbf 0}{}
\end{componenttable}

\begin{componenttable}[tab:components-zrevz]{Z^{\mathrm R} Z}
\componentrow{real scalar}{a^{2}+b^{2}+\mathbf A^{2}+\mathbf B^{2}}{}
\componentrow{imaginary scalar}{0}{}
\componentrow{real vector}{2(a\mathbf A+b\mathbf B-\mathbf A\times\mathbf B)}{\(\leftarrow\) differs from \(ZZ^{\mathrm R}\)}
\componentrow{imaginary vector}{\mathbf 0}{}
\end{componenttable}

\noindent The following two tables show that the conjugate-reverse products collapse to purely scalar and imaginary-scalar parts.
\begin{componenttable}[tab:components-zzcrev]{Z Z^{\mathrm{R}*}}
\componentrow{real scalar}{a^{2}-b^{2}-\mathbf A^{2}+\mathbf B^{2}}{}
\componentrow{imaginary scalar}{2ab-2\mathbf A\cdot\mathbf B}{}
\componentrow{real vector}{\mathbf 0}{}
\componentrow{imaginary vector}{\mathbf 0}{}
\end{componenttable}

\begin{componenttable}[tab:components-zcrevz]{Z^{\mathrm{R}*} Z}
\componentrow{real scalar}{a^{2}-b^{2}-\mathbf A^{2}+\mathbf B^{2}}{}
\componentrow{imaginary scalar}{2ab-2\mathbf A\cdot\mathbf B}{}
\componentrow{real vector}{\mathbf 0}{}
\componentrow{imaginary vector}{\mathbf 0}{}
\end{componenttable}
